\subsubsection{Impedance Analysis and Material Sensitivity}

The data reveals a pronounced inverse relationship between NaCl concentration and the magnitude of electrical impedance ($|Z|$) across the tested frequency spectrum. As the NaCl concentration increases, the impedance decreases, indicating an increase in conductivity. This trend is consistent with the expected behavior of ionic solutions, where higher salt content facilitates charge transport. We also observe a decrease in the magnitude of the impedance as the frequency increases. 

However, the addition of NaCl and increase of conductivity introduces a complex trade-off regarding the phase response. While increasing the NaCl concentration to 5\% effectively lowers the resistive component, the data indicates that it simultaneously induces a more capacitive response compared to the 0\% (or low-concentration) baseline.

Specifically, the phase angle ($\theta$) for the solution without NaCl remains closer to 0 degrees, indicating a predominantly resistive behavior. In contrast, the 5\% formulation, while lower than the 2.5\% solution, exhibits a larger phase shift. This increased capacitive reactance is attributable to interfacial polarization. The higher density of free ions facilitates the formation of a stronger electrical double-layer (EDL) at the electrode-gelatine interface, thereby increasing the effective capacitance of the system alongside its conductivity, especially at lower frequencies.

These findings align with the results obtained by Owda \emph{et al.}~\cite{Owda2020} in their study of gelatine-based phantoms, where they also reported similar trends in impedance and phase behavior with varying salt concentrations.

The reproducibility analysis between two batches of the 2.5\% NaCl gelatine samples demonstrated a high degree of consistency in both impedance magnitude and phase response. The negligible variance observed between batches indicates that the manufacturing process is reliable and produces uniform material properties, which is crucial for phantom applications requiring repeatable electrical characteristics.

\subsubsection{Signal Localization}

The spatial distribution of the measured voltages confirms the effective propagation and localization of the dipole signal within the conductive medium, as getting closer to each source results in a higher measured potential at the corresponding electrode frequency.
